\section{Blocks with dihedral defect groups: preliminary results}\label{sec:blocksWithDihedralDefectGroups}

To start with, we state two results, which do not require the defect of the blocks considered to be maximal.
First, we restate \cite[Lemma~4.1]{MacGregor22} in terms of splendid Morita equivalences.

\begin{lem}\label{lem:coveringOddIndex}%
    Let $G$ be a finite group and let $N\unlhd G$ be a normal subgroup such that $|G:N|$ is odd.  If $\bb\in\Bl(kN)$ is a block with dihedral defect groups of order $2^n$ $(n\geq 3)$ and $\bB\in \Bl(kG\mid \bb)$, then the $(\bB,\bb)$-bimodule $e_{\bB}kG e_{\bb}$ and its dual realize a splendid Morita equivalence between $\bB$ and $\bb$. (The blocks $\bb$ and $\bB$ are \emph{isomorphic via restriction} in the terminology of \cite[Theorem~4.1 and Definition~4.3]{HidaKoshitani}.)
\end{lem}

\begin{proof}
    First, observe that \cite[Lemma~4.1]{MacGregor22} proves that restriction from $G$ to $N$ induces a height-preserving bijection between $\Irr(\bB)$ and $\Irr(\bb)$. Thus, by \cite[0.1 Th\'{e}or\`{e}me and 0.2 Th\'{e}or\`{e}me]{BroueIsom90} the $(\bB,\bb)$-bimodule $e_{\bB}kG e_{\bb}$ and its dual realize a Morita equivalence between $\bB$ and $\bb$. This equivalence is a splendid Morita equivalence, since $e_{\bB}kG e_{\bb}$ is a $p$-permutation bimodule. 
\end{proof}

%\noindent{
%\color{blue}
%{[}Comment: (1) Proof works for semi-dihedral defect groups, too.{]}
%}\\  

 \begin{lem}\label{lem: nilpotentblock}
 Let $G$ be a finite soluble group. 
 %with a dihedral Sylow $2$-subgroup of order $2^n$ $(n\geq 3)$. 
 Then, any nilpotent  block $\bB\in\Bl(kG)$ with a defect group $D$ isomorphic to a dihedral $2$-group of order $2^n$ $(n\geq 3)$  is splendidly Morita equivalent to $kD_{2^n}$.
 \end{lem}
% \todo[color=green]{
 %Warning: $\bD_k(SD_{2^n})\cong \mathbb{Z}/2\mathbb{Z}$ and $\bD_k(Q_{2^n})\cong \mathbb{Z}/2\mathbb{Z}\oplus \mathbb{Z}/4\mathbb{Z} $, so we obtain $2$, resp. $8$, possible source algebras in these cases.
% }



\begin{proof}
    Let $\bB\in\Bl(kG)$ be a nilpotent block and let $A$ be a source algebra of $\bB$. By Puig's characterisation of the source algebras of nilpotent blocks \cite{PuigNilpotentBlocks88} (see \cite[Theorem 8.11.5]{LinckelmannVol2}), there exists an endo-permutation $kD$-module $V$ such that 
    $A\cong \End_k(V)\otimes_k kD$ as interior $D$-algebras. Moreover, $D$ is a vertex of $V$ and $V$ is a source of the unique simple $\bB$-module. Thus, $V$ is the trivial $kD$-module by Lemma~\ref{lem:sourcesSimpleDihedral}(b)  and the claim follows. 
 \end{proof}


 
%\noindent 
%????????\\
%????????\\
%????????\\
